\documentclass[10pt,a4paper]{article}

\usepackage[margin=1in]{geometry}

% Palatino font
\usepackage{newpxtext}
\usepackage{amssymb}
\usepackage{tikz}
\usetikzlibrary{positioning,shapes.multipart,arrows.meta, calc}
\usepackage{url}
\usepackage{graphicx}
\usepackage{amsmath}

\title{Neural Networks from Scratch and Exploring Compression}
\author{Robert Davidson}
\date{\today}

\begin{document}

\maketitle

\section{Feed Forward Neural Networks}

\subsection{Motivation}
Given some unknown function, $f^\ast: X \to Y$, we aim to find a good
approximation, $f: X \to Y$, by learning a set of parameters, $\theta$,
using a dataset of $N$ samples:
$$
\mathcal{D} = \{(\mathbf{x}_i, \mathbf{y}_i = f^\ast(\mathbf{x}_i))\}_{i=1}^N
$$

\noindent We measure how well our approximation performs by measuring the loss,
$L(\mathbf{y}_i, \hat{\mathbf{y}}_i)$. We then optimize the parameters, $\theta$, by
minimizing the loss over the dataset.\cite{goodfellow2016deep}\cite{smets2024mathematics}

\subsection{Neurons and Layers}

\noindent
A neural network is a function, $f : \mathbb{R}^n \to \mathbb{R}^m$ takes an input vector $\mathbf{x} \in \mathbb{R}^n$ and produces an output vector $\hat{\mathbf{y}} \in \mathbb{R}^m$ by passing the input through a series of layers of neurons.  \\[1ex]
\noindent The neuron is the basic building block of a neural network. When we pass a vector to the network, the neuron computes a weighted sum of the inputs and adds a bias term,
$$a = w_1 x_1 + w_2 x_2 + \ldots + w_n x_n + b = \mathbf{w}^T \mathbf{x} + b$$
This is commonly called the preactivation value. We then apply a non linear activation function, $\sigma$ to obtain the output of the neuron:
$$h = \sigma(a)$$
We commonly taken $\sigma = \text{ReLU}$, where $\text{ReLU}(x) = \max(0, x)$. \\
We say a layer is  a series of neurons that take the same input, apply their own unique weights and biases, and product a vector of outputs. For example, the first hidden layer of a neural network can be written:
$$\mathbf{h}^{(1)} = \begin{bmatrix}
    h^{(1)}_1 \\
    h^{(1)}_2 \\
    \vdots \\
    h^{(1)}_m 
\end{bmatrix}
= 
\begin{bmatrix}
    \sigma(w^{(1)}_{1,1} x_1 + w^{(1)}_{1,2} x_2 + \ldots + w^{(1)}_{1,n} x_n + b^{(1)}_1) \\
    \sigma(w^{(1)}_{2,1} x_1 + w^{(1)}_{2,2} x_2 + \ldots + w^{(1)}_{2,n} x_n + b^{(1)}_2) \\
    \vdots \\
    \sigma(w^{(1)}_{m,1} x_1 + w^{(1)}_{m,2} x_2 + \ldots + w^{(1)}_{m,n} x_n + b^{(1)}_m)
\end{bmatrix}
= 
 \sigma \left(
\begin{bmatrix}
    (w^{(1)})^T\\
    (w^{(1)})^T\\
    \vdots \\
    (w^{(1)})^T
\end{bmatrix}
  \right)
  = 
\sigma(W^{(1)}\mathbf{x} + \mathbf{b})
$$
The second hidden layer takes this vector as input and produces another vector of outputs:
$$\mathbf{h}^{(2)} = \sigma(W^{(2)}\mathbf{h}^{(1)} + \mathbf{b}^{(2)})$$
In general, the $l$-th layer of a neural network can be written as:
$$\mathbf{h}^{(l)} = \sigma(W^{(l)}\mathbf{h}^{(l-1)} + \mathbf{b}^{(l)})$$
\subsection{Questions}
We might ask ourselves, why do we need to use a non linear activation function? How do we know our choice of weights is correct?

\section{MNIST and Our Neural Network}
Mnist is a dataset of $28 \times 28$ pixel images of handwritten digits. Each image is labeled with the digits it represents from $0$ to $9$. Let $\mathcal{D}$ denote the dataset, then:
$$\mathcal{D} = \{(x_i, y_i)\}_{i=1}^{60,000}$$
where $x_i \in \mathbb{R}^{28 \times 28}$ and $y_i \in \{0, 1, 2, 3, 4, 5, 6, 7, 8, 9\}$. 

\noindent We use one hidden layer, with $256$ neurons, and a ReLU activation function. Our network can be written as:
$$\hat{y} = f(x) = \text{softmax}(W^{(2)}(\text{ReLU}(W^{(1)}x + b^{(1)}) + b^{(2)}))$$
with:
$$W_1 \in \mathbb{R}^{256 \times 784}, b_1 \in \mathbb{R}^{256}, W_2 \in \mathbb{R}^{10 \times 256}, b_2 \in \mathbb{R}^{10}$$
We initialize the weights and biases following Kaiming He's initialization scheme.

\pagebreak
\nocite{*}

\bibliographystyle{plain}
\bibliography{references}

\end{document}